\documentclass[10pt,a4paper]{article}
\usepackage[margin=0.6in]{geometry}
\usepackage{amsmath,amssymb,amsfonts,stmaryrd}
\usepackage{xcolor}
\usepackage{enumitem}
\usepackage{tcolorbox}
\usepackage{multicol}
\usepackage{listings}
\usepackage{hyperref}

\definecolor{isabelleblue}{RGB}{0, 70, 140}
\definecolor{isabellekw}{RGB}{120, 0, 120}
\definecolor{boxbg}{RGB}{248, 249, 250}

\tcbset{
  colback=boxbg,
  colframe=isabelleblue,
  fonttitle=\bfseries\large,
  boxrule=0.8pt,
  arc=3pt,
  left=8pt,
  right=8pt,
  top=6pt,
  bottom=6pt
}

\title{\textbf{\huge Proofs by Case Analysis in Isabelle}\\[0.3em]
\large Lecture Notes}
\author{\textbf{Course:} Computer-Assisted Theorem Proving (7CCSACTP)}
\date{}

\begin{document}

\maketitle
\vspace{-1.5em}

\begin{tcolorbox}[title={1. Context \& Intuition}]
\begin{multicols}{2}
\textbf{Revision of Previous Topics:}
\begin{itemize}[leftmargin=1.2em, itemsep=2pt]
    \item Isabelle interface, terms, and theorem statements.
    \item Forward reasoning (FW), Backward reasoning (BW).
    \item Proofs by substitution (\texttt{subst}).
    \item Mathematical induction for natural numbers (\texttt{nat}).
\end{itemize}

\columnbreak

\textbf{Case Analysis Intuition:}
\begin{itemize}[leftmargin=1.2em, itemsep=2pt]
    \item When proving a property for all numbers/lists, different cases may require distinct reasoning paths.
    \item \textbf{Example (Monus / Subtraction on \texttt{nat}):}
    \[ 3 - 1 = 2, \quad 3 - 2 = 1, \quad 3 - 3 = 0, \quad 3 - 4 = 0 \]
    In Isabelle, $x - y = 0$ whenever $x \le y$. Thus, proving statements with subtraction often requires a case split on $x \ge y$ vs. $x < y$.
\end{itemize}
\end{multicols}
\end{tcolorbox}

\vspace{0.3em}

\begin{tcolorbox}[title={2. Worked Example: Proving $x - y \le x$}]
\textbf{Theorem:} $x - y \le x \qquad (\text{Induction on } x)$

\vspace{0.3em}
\textbf{Base Case ($x = 0$):}
\begin{itemize}[leftmargin=1.2em, itemsep=2pt]
    \item $\textbf{goal}::$ $0 - y \le 0$
    \item \textbf{Proof steps:}
    \begin{enumerate}[leftmargin=1.2em, itemsep=1pt]
        \item $0 - y = 0$ \quad \textit{(by subst lem1: $0 - a = 0$)}
        \item $0 \le 0$ \quad \textit{(by lem2: $a \le a$)}
    \end{enumerate}
    \item \textsf{finally} $0 - y \le 0 \quad \checkmark$
\end{itemize}
\end{tcolorbox}

\vspace{0.3em}

\begin{tcolorbox}[title={3. Step Case ($x \to \text{Suc } x$) with Case Analysis}]
\textbf{Induction Hypothesis (I.H.):} $x - y \le x$ \\
\textbf{Goal:} $(\text{Suc } x) - y \le \text{Suc } x$

\vspace{0.3em}
\textbf{Proof by Case Analysis on $x \ge y$:}

\begin{multicols}{2}
\textbf{Case 1 ($x \ge y$):}
\begin{enumerate}[leftmargin=1.2em, itemsep=2pt]
    \item $(\text{Suc } x) - y = \text{Suc } (x - y)$ \quad \textit{(by lem1 [Case 1 assump])}
    \item $(\text{Suc } x) - y \le \text{Suc } (x - y)$ \quad \textit{(from 1, lem3)}
    \item $\dots \le \text{Suc } x$ \quad \textit{(using lem4 [OF I.H.]) $\checkmark$}
\end{enumerate}

\columnbreak

\textbf{Case 2 ($x \ngeq y$, i.e. $x < y$):}
\begin{enumerate}[leftmargin=1.2em, itemsep=2pt]
    \item $x < y$ \quad \textit{(from lem5 [OF Case 2 assump])}
    \item $\text{Suc } x \le y$ \quad \textit{(from 1)}
    \item $(\text{Suc } x) - y = 0$ \quad \textit{(from lem2 [OF 2])}
    \item $0 \le \text{Suc } x$ \quad \textit{(from lem6)}
    \item $(\text{Suc } x) - y \le \text{Suc } x$ \quad \textit{(using 4, subst 3) $\checkmark$}
\end{enumerate}
\end{multicols}
\end{tcolorbox}

\vspace{0.3em}

\begin{tcolorbox}[title={Auxiliary Lemmas Used in Proof}]
\begin{multicols}{2}
\begin{itemize}[leftmargin=1.2em, itemsep=2pt]
    \item \texttt{lem1:} $a \ge b \Longrightarrow (\text{Suc } a) - b = \text{Suc } (a - b)$
    \item \texttt{lem2:} $a \le b \Longrightarrow a - b = 0$
    \item \texttt{lem3:} $a \le a$
\end{itemize}
\columnbreak
\begin{itemize}[leftmargin=1.2em, itemsep=2pt]
    \item \texttt{lem4:} $a \le b \Longrightarrow \text{Suc } a \le \text{Suc } b$
    \item \texttt{lem5:} $a \ngeq b \Longrightarrow a < b$
    \item \texttt{lem6:} $0 \le a$
\end{itemize}
\end{multicols}
\end{tcolorbox}

\end{document}
